\documentclass{article}
\usepackage[utf8]{inputenc}
\usepackage{geometry}
\geometry{a4paper, margin=1in}
\usepackage{minted}

\title{UAV Software Engineering: Build Systems}
\author{CMake and Out-of-Source Builds}
\date{}

\begin{document}

\maketitle

\section{Theoretical Foundation: The Need for Build Systems}
In large-scale robotics projects like PX4 firmware or ROS 2 workspaces, software is composed of thousands of source files, 
external libraries, and cross-dependencies. Compiling these ecosystems manually via terminal commands (e.g., \texttt{g++ main.cpp -o node}) 
is unmaintainable. 

To solve this, the industry standard is \textbf{CMake}. 
CMake is not a compiler itself; rather, it is a cross-platform build system generator. 
It reads high-level configuration rules and automatically generates the low-level build files (such as Makefiles) 
required to compile the project efficiently.

\section{UAV System Application: Configuring \texttt{CMakeLists.txt}}
The core of CMake configuration is a text file strictly named \texttt{CMakeLists.txt}. 
To build a multi-threaded telemetry node, the minimal configuration requires defining the project, adding the executable, and linking necessary libraries:

\begin{itemize}
    \item \textbf{Basic Setup}: \texttt{cmake\_minimum\_required(VERSION 3.10)} and \texttt{project(telemetry\_node)} define the baseline.
    \item \textbf{Compilation Target}: \texttt{add\_executable(telemetry\_node main.cpp)} instructs the system to compile the source code into an executable file.
    \item \textbf{Dependency Linking}: While C++ Standard Template Library (STL) headers like \texttt{<queue>} or \texttt{<mutex>} 
    are natively included by the compiler, external OS-level libraries are not. To use hardware threads on Linux, we must explicitly 
    find and link the system's thread library using \texttt{find\_package(Threads REQUIRED)} and 
    \texttt{target\_link\_libraries(telemetry\_node Threads::Threads)}.
\end{itemize}

\section{Implementation Details: Out-of-Source Builds}
In professional C++ development, build artifacts (Makefiles, object files, binaries, and cache files) must never pollute the source code directory. 
This is solved by utilizing an \textbf{out-of-source build} strategy:

\begin{enumerate}
    \item Create an isolated directory: \texttt{mkdir build} and \texttt{cd build}.
    \item Generate build files targeting the parent directory: \texttt{cmake ..}
    \item Compile the binary: \texttt{make}
    \item Execute the resulting program locally: \texttt{./telemetry\_node}
\end{enumerate}

To maintain a clean version control history in Git, the \texttt{build/} directory must be explicitly added to the \texttt{.gitignore} file. 
This ensures that machine-specific compiled binaries and temporary artifacts are never committed to the remote repository.

\end{document}